\section{Syntaktická analýza}\label{sec:syntakticka-analyza}

Lexer už zdrojový text rozdělil na tokeny. Nyní potřebujeme rozhodnout,~zda jsou seřazeny podle pravidel programovacího jazyka,~a~zachytit jejich strukturu.

\emph{Syntaktický analyzátor} neboli \emph{parser} dostane posloupnost tokenů a~podle gramatiky rozhodne,~zda tvoří syntakticky správný program. Jeho výstupem bývá derivační nebo syntaktický strom,~případně zpráva o~syntaktické chybě.

\subsection{Bezkontextové gramatiky}

Obecná gramatika zavedená v~sekci \smartref{\showrefs}{sec:gramatiky}{o~gramatikách} dovoluje mít na levé straně pravidla celý řetězec. Syntaxi programovacích jazyků však obvykle popisujeme jednodušším druhem gramatik.

\begin{definition}[Bezkontextová gramatika]
    Gramatika $G=(N,T,P,S)$ je \emph{bezkontextová},~pokud má každé pravidlo tvar
    \[
        A\to\alpha,
    \]
    kde $A\in N$ je jediný neterminál a~$\alpha\in(N\cup T)^*$.
\end{definition}

Název vyjadřuje,~že neterminál $A$ smíme přepsat bez ohledu na symboly stojící kolem něj. Regulární gramatiky jsou zvláštním případem bezkontextových gramatik,~ale opačná inkluze neplatí. Například pravidlo $S\to aSb$ je bezkontextové,~nikoliv však regulární.

\subsection{Aritmetické výrazy}

Prioritu běžných aritmetických operátorů zachycuje gramatika
\[
\begin{aligned}
    E&\to E+T\mid E-T\mid T,\\
    T&\to T*F\mid T/F\mid F,\\
    F&\to (E)\mid\texttt{id}\mid\texttt{num}.
\end{aligned}
\]
Neterminál $E$ představuje výraz,~$T$ term a~$F$ faktor. Násobení a~dělení se objeví hlouběji ve stromu než sčítání a~odčítání,~a~proto se vyhodnotí dříve. Strom na obrázku \ref{fig:strom-vyrazu} odpovídá posloupnosti tokenů
\[
    \texttt{id}+\texttt{id}*\texttt{num}.
\]

\begin{figure}[h]
    \centering
    \begin{tikzpicture}[
    nt/.style={state,fill=orange!12,font=\small},
    term/.style={task,fill=green!10,minimum width=10mm,minimum height=7mm,font=\small}
]
    \node[nt] (e0) at (0,0) {$E$};
    \node[nt] (e1) at (-2.6,-1.4) {$E$};
    \node[term] (plus) at (0,-1.4) {$+$};
    \node[nt] (t1) at (2.6,-1.4) {$T$};

    \node[nt] (t0) at (-2.6,-2.8) {$T$};
    \node[nt] (t2) at (1.4,-2.8) {$T$};
    \node[term] (mul) at (2.6,-2.8) {$*$};
    \node[nt] (f2) at (3.8,-2.8) {$F$};

    \node[nt] (f0) at (-2.6,-4.2) {$F$};
    \node[nt] (f1) at (1.4,-4.2) {$F$};
    \node[term] (id1) at (-2.6,-5.5) {\texttt{id}};
    \node[term] (id2) at (1.4,-5.5) {\texttt{id}};
    \node[term] (num) at (3.8,-4.2) {\texttt{num}};

    \foreach \x/\y in {e0/e1,e0/plus,e0/t1,e1/t0,t1/t2,t1/mul,
        t1/f2,t0/f0,t2/f1,f0/id1,f1/id2,f2/num}
        \draw[thick] (\x)--(\y);
\end{tikzpicture}

    \caption{Derivační strom aritmetického výrazu.}
    \label{fig:strom-vyrazu}
\end{figure}

Derivační strom,~se kterým jsme se poprvé setkali v~sekci \smartref{\showrefs}{sec:gramatiky}{o~gramatikách},~obsahuje i~neterminály a~pomocné kroky gramatiky. Překladač si často vytvoří úspornější \emph{abstraktní syntaktický strom} (zkráceně AST),~v~němž zůstanou jen informace potřebné pro další zpracování. Ve stromu výrazu $2+3*4$ bude kořenem sčítání,~jeho levým synem číslo $2$ a~pravým synem násobení čísel $3$ a~$4$.

\subsection{Analýza shora dolů a~zdola nahoru}

Parser \emph{shora dolů} začíná počátečním neterminálem a~volí pravidla tak,~aby vytvořil vstupní posloupnost tokenů. Strom tedy staví od kořene k~listům. Jednoduchým příkladem je rekurzivní sestup,~v~němž každému neterminálu odpovídá jedna funkce.

Parser \emph{zdola nahoru} naopak začíná tokeny. Hledá úseky odpovídající pravým stranám pravidel a~nahrazuje je jejich levou stranou. Strom staví od listů ke kořeni. Praktické generátory parserů často používají některou z~variant této metody.

Levá rekurze v~pravidle $E\to E+T$ by pro přímočarý rekurzivní sestup vedla k~nekonečnému volání funkce pro $E$. Pro implementaci proto použijeme ekvivalentní tvar
\[
\begin{aligned}
    E&\to TE',&
    E'&\to +TE'\mid-TE'\mid\lambda,\\
    T&\to FT',&
    T'&\to *FT'\mid/FT'\mid\lambda,\\
    F&\to (E)\mid\texttt{id}\mid\texttt{num}.
\end{aligned}
\]
Opakování v~$E'$ a~$T'$ lze v~Pythonu přirozeně zapsat cyklem.

\subsection{Gramatika malého jazyka}

Celý průběžný jazyk popíšeme pravidly
\[
\begin{aligned}
    \mathit{Program}&\to\text{\textit{Příkaz}}\ \mathit{Program}\mid\lambda,\\
    \text{\textit{Příkaz}}&\to\texttt{let id = }E\texttt{;}
        \mid\texttt{print }E\texttt{;},
\end{aligned}
\]
doplněnými předchozími pravidly pro $E$,~$T$ a~$F$. Jazyk tedy dovoluje deklarovat číselnou proměnnou a~vypsat hodnotu výrazu.

Datové třídy v~programu \ref{lst:ast} popisují AST. Společní předkové \texttt{Expr} a~\texttt{Stmt} umožňují snadno rozlišovat výrazy a~příkazy.

\lstinputlisting[
    caption={Vrcholy abstraktního syntaktického stromu.},
    label={lst:ast}
]{04-kompilatory/code/mini_compiler/ast_nodes.py}

Parser v~programu \ref{lst:parser} přímo odpovídá gramatice. Metoda \texttt{\_statement} rozpoznává oba druhy příkazů. Metody \texttt{\_expression} a~\texttt{\_term} zpracovávají operátory,~zatímco \texttt{\_factor} číslo,~identifikátor nebo závorku.

\lstinputlisting[
    caption={Parser metodou rekurzivního sestupu.},
    label={lst:parser}
]{04-kompilatory/code/mini_compiler/parser.py}

Metoda \texttt{\_consume} zároveň kontroluje povinný následující token. Chybí-li například středník,~parser skončí s~chybou obsahující pozici neočekávaného tokenu. Kvalitní praktický parser se navíc pokusí najít vhodné místo,~od něhož může pokračovat a~oznámit více chyb během jednoho překladu.
